\section{G1: Emergence and Grokking}

\begin{frame}{Grokking as Delayed Generalization}
    \begin{columns}[T,onlytextwidth]
        \column{0.56\textwidth}
        \begin{itemize}
            \item Grokking: test performance improves after training performance has already converged.
            \item Proposal claim: controlled train/test distribution shift is a strong inducer of this transition.
            \item Focus metric: transition delay and post-transition stability.
        \end{itemize}
        \footnotesize
        Key reference: \cite{carvalho2025inducing_grokking_distribution_shifts}
        \column{0.44\textwidth}
        \begin{figure}
            \centering
            \includegraphics[width=\linewidth]{grokking_definition.pdf}
            \caption{Definition and phase regions used in this proposal.}
        \end{figure}
    \end{columns}
\end{frame}

\begin{frame}{Synthetic Datasets Used for Controlled Shift}
    \begin{columns}[T,onlytextwidth]
        \column{0.49\textwidth}
        \centering
        \includegraphics[width=0.8\linewidth]{group_eqdist.pdf}
        
        {\scriptsize Equidistant subclasses}
        \column{0.49\textwidth}
        \centering
        \includegraphics[width=0.64\linewidth]{group_eqvariant.png}
        
        {\scriptsize Equivariant subclasses}
    \end{columns}
    \small
    These datasets are designed to manipulate class/subclass coverage and isolate how shift structure impacts late generalization.
\end{frame}

\begin{frame}{Representative G1 Results}
    \begin{columns}[T,onlytextwidth]
        \column{0.49\textwidth}
        \begin{figure}
            \centering
            \includegraphics[width=\linewidth]{result_small.pdf}
            \caption{Small imbalance can induce late generalization.}
        \end{figure}
        \column{0.49\textwidth}
        \begin{figure}
            \centering
            \includegraphics[width=\linewidth]{result_large.pdf}
            \caption{Shift-driven behavior remains visible at larger scale with strong imbalance.}
        \end{figure}
    \end{columns}
\end{frame}
